\documentclass[11pt,a4paper]{article}

% --- THIẾT LẬP LỀ TRANG & KHỔ GIẤY ---
\usepackage[top=1.8cm,bottom=1.8cm,left=1.8cm,right=1.8cm]{geometry}

% --- HỖ TRỢ ĐA DẠNG TRÌNH BIÊN DỊCH (XeLaTeX / pdfLaTeX / LuaLaTeX) ---
\usepackage{iftex}
\ifXeTeX
  \usepackage{fontspec}
  \usepackage[vietnamese]{babel}
  \defaultfontfeatures{Ligatures=TeX}
\else
  \usepackage[utf8]{inputenc}
  \usepackage[T5,T1]{fontenc}
  \usepackage[vietnamese]{babel}
  \usepackage{mathptmx}
\fi

% --- CÁC GÓI TOÁN HỌC CHUẨN AMS ---
\usepackage{amsmath}
\usepackage{amssymb}
\usepackage{amsfonts}
\usepackage{mathtools}
\usepackage{bm}

% --- BẢNG BIỂU, DANH SÁCH & TRANG TRÍ ---
\usepackage{tabularx}
\usepackage{booktabs}
\usepackage{array}
\usepackage{enumitem}
\usepackage{multicol}
\usepackage{fancyhdr}
\usepackage{xcolor}

% --- THIẾT LẬP HEADER & FOOTER ---
\pagestyle{fancy}
\fancyhf{}
\renewcommand{\headrulewidth}{0pt}
\lfoot{\textit{\small Bộ đề kiểm tra Toán THCS KNTT - Mã: c1\_6\_bai5\_test2}}
\rfoot{\textit{\small Trang \thepage}}

\begin{document}

% ==============================================================================
% PHẦN ĐẦU ĐỀ THI (HEADER TIÊU CHUẨN BỘ GD&ĐT)
% ==============================================================================
\noindent
\begin{tabular*}{\textwidth}{@{\extracolsep{\fill}}lr@{}}
\textbf{SỞ GD\&ĐT ........................} & \textbf{ĐỀ KIỂM TRA ĐỊNH KỲ NĂM HỌC 2025 - 2026} \\
\textbf{TRƯỜNG THCS ........................} & \textbf{Môn: TOÁN - LỚP 6} \\
\textit{(Đề thi chuẩn ma trận 2025)} & \textit{Thời gian làm bài: 45 phút (không kể giao đề)}
\end{tabular*}

\vspace{6pt}
\begin{center}
{\Large\bfseries Toán 6 (KNTT) - Bài 5: Phép nhân và phép chia số tự nhiên - Bài kiểm tra số 2}\\[4pt]
{\bfseries [ MỨC ĐỘ 2: KHÁ - Mã đề: \texttt{c1\_6\_bai5\_test2} ]}
\end{center}

\vspace{4pt}
\noindent
\rule{\textwidth}{0.8pt}\\[3pt]
Họ và tên thí sinh: \dotfill\ Lớp: \dots\dots\ SBD: \dots\dots\\[-6pt]
\rule{\textwidth}{0.8pt}

\vspace{10pt}

% ==============================================================================
% PHẦN I: TRẮC NGHIỆM NHIỀU PHƯƠNG ÁN (3,0 ĐIỂM)
% ==============================================================================
\noindent
\textbf{\large PHẦN I. Câu trắc nghiệm nhiều phương án lựa chọn (3,0 điểm)}\\
\textit{\small Thí sinh trả lời từ câu 1 đến câu 12. Mỗi câu hỏi thí sinh chỉ chọn một phương án đúng nhất.}
\vspace{6pt}
\noindent\textbf{Câu 1.} [Nhận biết] Trong một phép chia có dư với số chia bằng$18$, số dư lớn nhất có thể có của phép chia đó là:\par\vspace{2pt}
\noindent\begin{tabular*}{\textwidth}{@{\extracolsep{\fill}}llll@{}}
\textbf{A.} $17$ & \textbf{B.} $18$ & \textbf{C.} $19$ & \textbf{D.} $0$
\end{tabular*}\par\vspace{5pt}

\noindent\textbf{Câu 2.} [Nhận biết] Cho tích$A = 25 \cdot 48 \cdot 4$. Áp dụng tính chất giao hoán và kết hợp, cách viết nào dưới đây giúp tính nhẩm nhanh nhất?\par\vspace{2pt}
\noindent\begin{tabular*}{\textwidth}{@{\extracolsep{\fill}}p{0.48\textwidth}p{0.48\textwidth}@{}}
\textbf{A.} $(25 \cdot 4) \cdot 48$ & \textbf{B.} $(25 \cdot 48) \cdot 4$ \\
\textbf{C.} $25 \cdot (48 + 4)$ & \textbf{D.} $(25 + 4) \cdot 48$
\end{tabular*}\par\vspace{5pt}

\noindent\textbf{Câu 3.} [Nhận biết] Cho$a, b \in \mathbb{N}$. Biết rằng$a \cdot b = 0$, khẳng định nào sau đây chắc chắn đúng?\par\vspace{2pt}
\noindent\begin{tabular*}{\textwidth}{@{\extracolsep{\fill}}p{0.48\textwidth}p{0.48\textwidth}@{}}
\textbf{A.} $a = 0$và$b = 0$ & \textbf{B.} $a = 0$hoặc$b = 0$ \\
\textbf{C.} $a \ne 0$và$b = 0$ & \textbf{D.} $a + b = 1$
\end{tabular*}\par\vspace{5pt}

\noindent\textbf{Câu 4.} [Nhận biết] Khi thực hiện phép chia có dư số tự nhiên$a$cho số tự nhiên$b$($b \ne 0$) được thương là$q$và số dư là$r$. Khẳng định đúng về số dư$r$là:\par\vspace{2pt}
\noindent\begin{tabular*}{\textwidth}{@{\extracolsep{\fill}}llll@{}}
\textbf{A.} $0 < r < b$ & \textbf{B.} $0 \le r \le b$ & \textbf{C.} $0 \le r < b$ & \textbf{D.} $0 < r \le b$
\end{tabular*}\par\vspace{5pt}

\noindent\textbf{Câu 5.} [Thông hiểu] Giá trị của biểu thức$M = 37 \cdot 64 + 37 \cdot 36 - 700$bằng:\par\vspace{2pt}
\noindent\begin{tabular*}{\textwidth}{@{\extracolsep{\fill}}llll@{}}
\textbf{A.} $3\,000$ & \textbf{B.} $3\,700$ & \textbf{C.} $2\,700$ & \textbf{D.} $3\,100$
\end{tabular*}\par\vspace{5pt}

\noindent\textbf{Câu 6.} [Thông hiểu] Số tự nhiên$x$thỏa mãn$125 - 5 \cdot (x + 4) = 35$là:\par\vspace{2pt}
\noindent\begin{tabular*}{\textwidth}{@{\extracolsep{\fill}}llll@{}}
\textbf{A.} $14$ & \textbf{B.} $18$ & \textbf{C.} $22$ & \textbf{D.} $10$
\end{tabular*}\par\vspace{5pt}

\noindent\textbf{Câu 7.} [Thông hiểu] Một hội trường cần sắp xếp chỗ ngồi cho$325$đại biểu. Biết mỗi dãy ghế dài có đúng$16$chỗ ngồi. Hỏi ban tổ chức cần ít nhất bao nhiêu dãy ghế như vậy để tất cả đại biểu đều có chỗ ngồi?\par\vspace{2pt}
\noindent\begin{tabular*}{\textwidth}{@{\extracolsep{\fill}}llll@{}}
\textbf{A.} $20$dãy ghế & \textbf{B.} $21$dãy ghế & \textbf{C.} $22$dãy ghế & \textbf{D.} $19$dãy ghế
\end{tabular*}\par\vspace{5pt}

\noindent\textbf{Câu 8.} [Thông hiểu] Khi chia một số tự nhiên cho$42$ta được số thương là$15$và số dư là$28$. Nếu đem số tự nhiên đó chia cho$14$thì kết luận nào sau đây là đúng?\par\vspace{2pt}
\noindent\textbf{A.} Phép chia hết và có số dư bằng$0$\par\vspace{1pt}
\noindent\textbf{B.} Phép chia có số dư là$7$\par\vspace{1pt}
\noindent\textbf{C.} Phép chia có số dư là$14$\par\vspace{1pt}
\noindent\textbf{D.} Phép chia có thương là$45$và số dư là$2$\par\vspace{5pt}

\noindent\textbf{Câu 9.} [Thông hiểu] Tích của hai số tự nhiên bằng$864$. Nếu một thừa số tăng lên gấp$3$lần và thừa số còn lại giảm đi$2$lần thì tích mới thu được là:\par\vspace{2pt}
\noindent\begin{tabular*}{\textwidth}{@{\extracolsep{\fill}}llll@{}}
\textbf{A.} $1\,296$ & \textbf{B.} $576$ & \textbf{C.} $2\,592$ & \textbf{D.} $432$
\end{tabular*}\par\vspace{5pt}

\noindent\textbf{Câu 10.} [Thông hiểu] Tính nhanh giá trị của biểu thức$B = 234 \cdot 101 - 234$ta được kết quả là:\par\vspace{2pt}
\noindent\begin{tabular*}{\textwidth}{@{\extracolsep{\fill}}llll@{}}
\textbf{A.} $23\,400$ & \textbf{B.} $23\,634$ & \textbf{C.} $2\,340$ & \textbf{D.} $23\,166$
\end{tabular*}\par\vspace{5pt}

\noindent\textbf{Câu 11.} [Thông hiểu] Một đoàn khách gồm$185$người cần sang sông bằng thuyền nhỏ. Mỗi chuyến thuyền chỉ chở tối đa$12$khách (không kể người lái thuyền). Hỏi người lái thuyền cần thực hiện ít nhất bao nhiêu chuyến để chở hết số khách qua sông?\par\vspace{2pt}
\noindent\begin{tabular*}{\textwidth}{@{\extracolsep{\fill}}llll@{}}
\textbf{A.} $15$chuyến & \textbf{B.} $16$chuyến & \textbf{C.} $17$chuyến & \textbf{D.} $14$chuyến
\end{tabular*}\par\vspace{5pt}

\noindent\textbf{Câu 12.} [Thông hiểu] Tích$P = 1 \cdot 2 \cdot 3 \cdot 4 \dots \cdot 15$tận cùng bằng bao nhiêu chữ số$0$liên tiếp?\par\vspace{2pt}
\noindent\begin{tabular*}{\textwidth}{@{\extracolsep{\fill}}llll@{}}
\textbf{A.} $3$ & \textbf{B.} $2$ & \textbf{C.} $4$ & \textbf{D.} $1$
\end{tabular*}\par\vspace{5pt}


% ==============================================================================
% PHẦN II: TRẮC NGHIỆM ĐÚNG / SAI (2,0 ĐIỂM)
% ==============================================================================
\vspace{6pt}
\noindent
\textbf{\large PHẦN II. Câu trắc nghiệm đúng sai (2,0 điểm)}\\
\textit{\small Thí sinh trả lời từ câu 1 đến câu 2. Trong mỗi ý a), b), c), d) ở mỗi câu, thí sinh chọn đúng hoặc sai.}\\
\textit{\small (Quy tắc tính điểm: Đúng 1 ý được 0,1đ; đúng 2 ý được 0,25đ; đúng 3 ý được 0,5đ; đúng 4 ý được 1,0đ).}
\vspace{6pt}
\noindent\textbf{Câu 1.} Cho hai biểu thức số$A = 125 \cdot 72$và$B = 48 \cdot 15 + 48 \cdot 85$.\par\vspace{2pt}
\begin{enumerate}[label=\textbf{\alph*)}, leftmargin=1.2cm, itemsep=2pt, topsep=2pt]
  \item \textnormal{[Thông hiểu]} Giá trị của biểu thức$B$có thể tính nhanh bằng cách đưa về$48 \cdot (15 + 85) = 4\,800$.
  \item \textnormal{[Thông hiểu]} Giá trị của biểu thức$A$bằng$9\,000$.
  \item \textnormal{[Vận dụng thấp]} Hiệu$A - B$chia hết cho cả ba số$2, 5$và$10$.
  \item \textnormal{[Vận dụng thấp]} Khi thực hiện phép chia có dư$A$cho$B$, ta được số thương là$2$và số dư là$600$.
\end{enumerate}
\vspace{6pt}

\noindent\textbf{Câu 2.} Học sinh lớp 6A gom được$468$quyển vở để đóng gói gửi tặng các bạn học sinh vùng cao khó khăn.\par\vspace{2pt}
\begin{enumerate}[label=\textbf{\alph*)}, leftmargin=1.2cm, itemsep=2pt, topsep=2pt]
  \item \textnormal{[Thông hiểu]} Nếu mỗi gói quà cần đóng đúng$15$quyển vở thì lớp 6A đóng được nhiều nhất$31$gói quà và còn dư$3$quyển vở.
  \item \textnormal{[Thông hiểu]} Để đóng thêm được đúng$1$gói quà loại$15$quyển vở nữa từ số vở còn thừa, lớp 6A phải quyên góp thêm ít nhất$12$quyển vở.
  \item \textnormal{[Vận dụng thấp]} Nếu mỗi gói quà giảm xuống chỉ đóng$12$quyển vở thì toàn bộ$468$quyển vở sẽ được chia đều mà không còn thừa quyển nào.
  \item \textnormal{[Vận dụng thấp]} Nếu đem toàn bộ số vở đó chia đều cho$18$học sinh nghèo thì mỗi bạn nhận được$25$quyển và còn dư$18$quyển.
\end{enumerate}
\vspace{6pt}


% ==============================================================================
% PHẦN III: TRẮC NGHIỆM TRẢ LỜI NGẮN (2,0 ĐIỂM)
% ==============================================================================
\vspace{6pt}
\noindent
\textbf{\large PHẦN III. Câu trắc nghiệm trả lời ngắn (2,0 điểm)}\\
\textit{\small Thí sinh trả lời từ câu 1 đến câu 4. Mỗi câu trả lời đúng thí sinh được 0,5 điểm.}
\vspace{6pt}
\noindent\textbf{Câu 1.} [Thông hiểu] Tìm số tự nhiên$x$thỏa mãn:$(3x - 15) : 7 = 12$.\par\vspace{2pt}
\noindent\textit{Đáp số:} \dotfill\par\vspace{6pt}

\noindent\textbf{Câu 2.} [Thông hiểu] Một số tự nhiên$N$khi chia cho$24$có số dư là$17$. Hỏi số$N$khi chia cho$8$có số dư là bao nhiêu?\par\vspace{2pt}
\noindent\textit{Đáp số:} \dotfill\par\vspace{6pt}

\noindent\textbf{Câu 3.} [Vận dụng thấp] Tính tổng của dãy số chẵn liên tiếp:$S = 2 + 4 + 6 + 8 + \dots + 100$.\par\vspace{2pt}
\noindent\textit{Đáp số:} \dotfill\par\vspace{6pt}

\noindent\textbf{Câu 4.} [Vận dụng thấp] Bác An mua$15$hộp sữa chua với giá$8$nghìn đồng một hộp và$8$gói bánh với giá$25$nghìn đồng một gói. Bác An đưa cho thu ngân tờ tiền mệnh giá$400$nghìn đồng. Hỏi bác An được nhận lại bao nhiêu nghìn đồng tiền thừa? (Chỉ điền số nguyên, đơn vị là nghìn đồng).\par\vspace{2pt}
\noindent\textit{Đáp số:} \dotfill\par\vspace{6pt}


% ==============================================================================
% PHẦN IV: TỰ LUẬN (3,0 ĐIỂM)
% ==============================================================================
\vspace{6pt}
\noindent
\textbf{\large PHẦN IV. Tự luận (3,0 điểm)}\\
\textit{\small Thí sinh trình bày chi tiết các bước giải từ câu 1 đến câu 3 vào giấy làm bài.}
\vspace{6pt}
\noindent\textbf{Bài 1 (1,0 điểm).} [Thông hiểu] Tính giá trị của biểu thức sau một cách hợp lý:$$A = 135 \cdot 48 + 52 \cdot 135 - 35 \cdot 100$$\par\vspace{8pt}

\noindent\textbf{Bài 2 (1,0 điểm).} [Vận dụng thấp] Một trường THCS tổ chức cho$345$học sinh khối 6 tham gia dã ngoại. Nhà trường dự kiến sử dụng hai loại ô tô: xe lớn loại$45$chỗ và xe nhỏ loại$29$chỗ (không tính ghế lái xe).\par\noindent a) Nhà trường đã thuê$6$xe loại$45$chỗ. Hỏi nhà trường cần thuê thêm ít nhất bao nhiêu xe loại$29$chỗ để đảm bảo tất cả các học sinh còn lại đều có chỗ ngồi?\par\noindent b) Biết giá thuê một xe$45$chỗ là$3\,200\,000$đồng và một xe$29$chỗ là$2\,100\,000$đồng. Tính tổng số tiền nhà trường phải chi trả để thuê toàn bộ số xe trên.\par\vspace{8pt}

\noindent\textbf{Bài 3 (1,0 điểm).} [Vận dụng cao] Tìm tất cả các số tự nhiên$n$sao cho biểu thức$(3n + 16)$chia hết cho$(n + 2)$.\par\vspace{8pt}


\begin{center}
\textbf{---------- HẾT ----------}\\
\textit{\small Cán bộ coi thi không giải thích gì thêm.}
\end{center}

\newpage
% ==============================================================================
% HƯỚNG DẪN CHẤM VÀ ĐÁP ÁN CHI TIẾT
% ==============================================================================
\begin{center}
{\Large\bfseries HƯỚNG DẪN CHẤM VÀ ĐÁP ÁN CHI TIẾT}\\[4pt]
{\large\bfseries Toán 6 (KNTT) - Bài 5: Phép nhân và phép chia số tự nhiên - Bài kiểm tra số 2}\\[2pt]
{\bfseries Mã đề: \texttt{c1\_6\_bai5\_test2}}
\end{center}

\vspace{10pt}
\noindent\textbf{\large 1. ĐÁP ÁN PHẦN I: TRẮC NGHIỆM NHIỀU PHƯƠNG ÁN (3,0 điểm)}
\vspace{4pt}
\begin{center}
\begin{tabular}{|c|c|c|c|c|c|c|c|c|c|c|c|}
\hline
\textbf{1} & \textbf{2} & \textbf{3} & \textbf{4} & \textbf{5} & \textbf{6} & \textbf{7} & \textbf{8} & \textbf{9} & \textbf{10} & \textbf{11} & \textbf{12} \\
\hline
\textbf{A} & \textbf{A} & \textbf{B} & \textbf{C} & \textbf{A} & \textbf{A} & \textbf{B} & \textbf{A} & \textbf{A} & \textbf{A} & \textbf{B} & \textbf{A} \\
\hline
\end{tabular}
\end{center}
\vspace{6pt}


\noindent\textbf{\large 2. ĐÁP ÁN PHẦN II: TRẮC NGHIỆM ĐÚNG / SAI (2,0 điểm)}
\vspace{4pt}
\begin{center}
\begin{tabular}{|c|c|c|c|c|}
\hline
\textbf{Câu hỏi} & \textbf{Ý a)} & \textbf{Ý b)} & \textbf{Ý c)} & \textbf{Ý d)} \\
\hline
\textbf{Câu 1} & \textbf{Đúng} & \textbf{Đúng} & \textbf{Đúng} & \textbf{Sai} \\
\hline
\textbf{Câu 2} & \textbf{Đúng} & \textbf{Đúng} & \textbf{Đúng} & \textbf{Sai} \\
\hline
\end{tabular}
\end{center}
\vspace{6pt}


\noindent\textbf{\large 3. ĐÁP ÁN PHẦN III: TRẢ LỜI NGẮN (2,0 điểm)}
\vspace{4pt}
\begin{center}
\begin{tabular}{|c|m{11cm}|c|}
\hline
\textbf{Câu} & \textbf{Tóm tắt yêu cầu} & \textbf{Đáp số chính xác} \\
\hline
\textbf{Câu 1} & Tìm số tự nhiên$x$thỏa mãn:$(3x - 15) : 7 = 12$. & \textbf{33} \\
\hline
\textbf{Câu 2} & Một số tự nhiên$N$khi chia cho$24$có số dư là$17$. Hỏi số$N$khi chia cho... & \textbf{1} \\
\hline
\textbf{Câu 3} & Tính tổng của dãy số chẵn liên tiếp:$S = 2 + 4 + 6 + 8 + \dots + 100$. & \textbf{2550} \\
\hline
\textbf{Câu 4} & Bác An mua$15$hộp sữa chua với giá$8$nghìn đồng một hộp và$8$gói bánh với... & \textbf{80} \\
\hline
\end{tabular}
\end{center}
\vspace{6pt}


\noindent\textbf{\large 4. LỜI GIẢI CHI TIẾT VÀ BIỂU ĐIỂM TỰ LUẬN (3,0 điểm)}
\vspace{4pt}
\noindent\textbf{\underline{Bài 1 (1,0 điểm):}}\par\vspace{3pt}
Bước 1 (0.5đ): Áp dụng tính chất phân phối của phép nhân đối với phép cộng:$A = 135 \cdot (48 + 52) - 35 \cdot 100 = 135 \cdot 100 - 35 \cdot 100$.\par\noindent Bước 2 (0.5đ): Tiếp tục đặt thừa số chung hoặc tính trực tiếp:$A = 100 \cdot (135 - 35) = 100 \cdot 100 = 10\,000$(hoặc$13\,500 - 3\,500 = 10\,000$).\par\vspace{8pt}

\noindent\textbf{\underline{Bài 2 (1,0 điểm):}}\par\vspace{3pt}
Bước 1 (0.5đ): Tính số chỗ ngồi của$6$xe$45$chỗ là:$6 \cdot 45 = 270$(chỗ). Số học sinh còn lại là:$345 - 270 = 75$(học sinh).\par\noindent Bước 2 (0.25đ): Thực hiện phép chia:$75 : 29 = 2$(dư 17). Vì còn dư$17$học sinh nên nhà trường phải thuê thêm ít nhất:$2 + 1 = 3$(xe$29$chỗ).\par\noindent Bước 3 (0.25đ): Tổng số tiền nhà trường phải trả là:$6 \cdot 3\,200\,000 + 3 \cdot 2\,100\,000 = 19\,200\,000 + 6\,300\,000 = 25\,500\,000$(đồng).\par\vspace{8pt}

\noindent\textbf{\underline{Bài 3 (1,0 điểm):}}\par\vspace{3pt}
Bước 1 (0.5đ): Ta biến đổi:$3n + 16 = 3(n + 2) + 10$. Vì$3(n + 2)$luôn chia hết cho$(n + 2)$với mọi số tự nhiên$n$, nên để$(3n + 16) \,\vdots\, (n + 2)$thì$10$phải chia hết cho$(n + 2)$.\par\noindent Bước 2 (0.5đ): Do đó$(n + 2)$phải là ước tự nhiên của$10$. Các ước tự nhiên của$10$là$1; 2; 5; 10$. Vì$n$là số tự nhiên nên$n + 2 \ge 2$. Do đó:\par\noindent -$n + 2 = 2 \Rightarrow n = 0$;\par\noindent -$n + 2 = 5 \Rightarrow n = 3$;\par\noindent -$n + 2 = 10 \Rightarrow n = 8$.\par\noindent Kết luận:$n \in \{0; 3; 8\}$.\par\vspace{8pt}

\end{document}
